\documentclass[11pt]{article}
% Preambulo comun para los informes de la ronda 3.
% Se incluye con % Preambulo comun para los informes de la ronda 3.
% Se incluye con % Preambulo comun para los informes de la ronda 3.
% Se incluye con \input{../comun/preambulo} desde cada subcarpeta.
\usepackage[utf8]{inputenc}
\usepackage[T1]{fontenc}
\usepackage[spanish,es-nodecimaldot]{babel}
\usepackage{lmodern}
\usepackage[a4paper,margin=2.4cm]{geometry}
\usepackage{booktabs}
\usepackage{longtable}
\usepackage{array}
\usepackage{xcolor}
\usepackage{hyperref}
\usepackage{enumitem}
\usepackage{fancyhdr}
\usepackage{titlesec}
\usepackage{graphicx}
\usepackage{amssymb}
\usepackage{microtype}

\definecolor{sosTeal}{HTML}{0F766E}
\definecolor{sosGray}{HTML}{475569}

\hypersetup{
  colorlinks=true,
  linkcolor=sosTeal,
  urlcolor=sosTeal,
  citecolor=sosTeal,
  pdftitle={SOS Clinic - Ronda 3 de pruebas},
  pdfauthor={Equipo SOS Clinic}
}

\pagestyle{fancy}
\fancyhf{}
\lhead{\textcolor{sosGray}{SOS Clinic --- Ronda 3}}
\rhead{\textcolor{sosGray}{\leftmark}}
\cfoot{\thepage}
\setlength{\headheight}{15pt}
\renewcommand{\headrulewidth}{0.4pt}

\titleformat{\section}{\Large\bfseries\color{sosTeal}}{\thesection}{0.8em}{}
\titleformat{\subsection}{\large\bfseries\color{sosGray}}{\thesubsection}{0.6em}{}

\setlist[itemize]{leftmargin=1.4em,itemsep=2pt,topsep=4pt}
\setlist[enumerate]{leftmargin=1.6em,itemsep=2pt,topsep=4pt}

\newcommand{\pantalla}[1]{\texttt{\detokenize{#1}}}
\newcommand{\ok}{\textcolor{sosTeal}{\textbf{PASS}}}
\newcommand{\fallo}{\textcolor{red!70!black}{\textbf{FAIL}}}
\newcommand{\aviso}{\textcolor{orange!80!black}{\textbf{WARN}}}
\newcommand{\omitido}{\textcolor{sosGray}{\textbf{SKIP}}}
 desde cada subcarpeta.
\usepackage[utf8]{inputenc}
\usepackage[T1]{fontenc}
\usepackage[spanish,es-nodecimaldot]{babel}
\usepackage{lmodern}
\usepackage[a4paper,margin=2.4cm]{geometry}
\usepackage{booktabs}
\usepackage{longtable}
\usepackage{array}
\usepackage{xcolor}
\usepackage{hyperref}
\usepackage{enumitem}
\usepackage{fancyhdr}
\usepackage{titlesec}
\usepackage{graphicx}
\usepackage{amssymb}
\usepackage{microtype}

\definecolor{sosTeal}{HTML}{0F766E}
\definecolor{sosGray}{HTML}{475569}

\hypersetup{
  colorlinks=true,
  linkcolor=sosTeal,
  urlcolor=sosTeal,
  citecolor=sosTeal,
  pdftitle={SOS Clinic - Ronda 3 de pruebas},
  pdfauthor={Equipo SOS Clinic}
}

\pagestyle{fancy}
\fancyhf{}
\lhead{\textcolor{sosGray}{SOS Clinic --- Ronda 3}}
\rhead{\textcolor{sosGray}{\leftmark}}
\cfoot{\thepage}
\setlength{\headheight}{15pt}
\renewcommand{\headrulewidth}{0.4pt}

\titleformat{\section}{\Large\bfseries\color{sosTeal}}{\thesection}{0.8em}{}
\titleformat{\subsection}{\large\bfseries\color{sosGray}}{\thesubsection}{0.6em}{}

\setlist[itemize]{leftmargin=1.4em,itemsep=2pt,topsep=4pt}
\setlist[enumerate]{leftmargin=1.6em,itemsep=2pt,topsep=4pt}

\newcommand{\pantalla}[1]{\texttt{\detokenize{#1}}}
\newcommand{\ok}{\textcolor{sosTeal}{\textbf{PASS}}}
\newcommand{\fallo}{\textcolor{red!70!black}{\textbf{FAIL}}}
\newcommand{\aviso}{\textcolor{orange!80!black}{\textbf{WARN}}}
\newcommand{\omitido}{\textcolor{sosGray}{\textbf{SKIP}}}
 desde cada subcarpeta.
\usepackage[utf8]{inputenc}
\usepackage[T1]{fontenc}
\usepackage[spanish,es-nodecimaldot]{babel}
\usepackage{lmodern}
\usepackage[a4paper,margin=2.4cm]{geometry}
\usepackage{booktabs}
\usepackage{longtable}
\usepackage{array}
\usepackage{xcolor}
\usepackage{hyperref}
\usepackage{enumitem}
\usepackage{fancyhdr}
\usepackage{titlesec}
\usepackage{graphicx}
\usepackage{amssymb}
\usepackage{microtype}

\definecolor{sosTeal}{HTML}{0F766E}
\definecolor{sosGray}{HTML}{475569}

\hypersetup{
  colorlinks=true,
  linkcolor=sosTeal,
  urlcolor=sosTeal,
  citecolor=sosTeal,
  pdftitle={SOS Clinic - Ronda 3 de pruebas},
  pdfauthor={Equipo SOS Clinic}
}

\pagestyle{fancy}
\fancyhf{}
\lhead{\textcolor{sosGray}{SOS Clinic --- Ronda 3}}
\rhead{\textcolor{sosGray}{\leftmark}}
\cfoot{\thepage}
\setlength{\headheight}{15pt}
\renewcommand{\headrulewidth}{0.4pt}

\titleformat{\section}{\Large\bfseries\color{sosTeal}}{\thesection}{0.8em}{}
\titleformat{\subsection}{\large\bfseries\color{sosGray}}{\thesubsection}{0.6em}{}

\setlist[itemize]{leftmargin=1.4em,itemsep=2pt,topsep=4pt}
\setlist[enumerate]{leftmargin=1.6em,itemsep=2pt,topsep=4pt}

\newcommand{\pantalla}[1]{\texttt{\detokenize{#1}}}
\newcommand{\ok}{\textcolor{sosTeal}{\textbf{PASS}}}
\newcommand{\fallo}{\textcolor{red!70!black}{\textbf{FAIL}}}
\newcommand{\aviso}{\textcolor{orange!80!black}{\textbf{WARN}}}
\newcommand{\omitido}{\textcolor{sosGray}{\textbf{SKIP}}}


\title{\textbf{Pruebas funcionales E2E}\\[0.3em]
\large Metodología y ejecución --- Ronda 3}
\author{Equipo SOS Clinic}
\date{22 de octubre de 2026}

\begin{document}
\maketitle

\begin{abstract}
Este documento describe cómo se planificaron y ejecutaron las pruebas
funcionales de punta a punta de SOS Clinic sobre el stack local completo:
autenticación con 2FA, gestión de pacientes, triage, signos vitales, alta con
notificaciones, analítica, agente conversacional, auditoría y turnos. El
objetivo es que cualquier integrante del equipo pueda reproducir la ejecución
sin depender de conocimiento tácito.
\end{abstract}

\tableofcontents

\section{Objetivo y alcance}

Se validó el comportamiento funcional observable por API a través del gateway
(\texttt{http://127.0.0.1:8080}) y el frontend estático
(\texttt{http://127.0.0.1:5500}). Quedaron fuera del alcance la automatización
completa de la interfaz con Playwright y las pruebas de carga; el frontend se
verificó con un smoke test de páginas y con la prueba de XSS en navegador real
que se documenta en el informe de pentest.

\section{Entorno de ejecución}

\begin{table}[h]
\centering
\small
\begin{tabular}{@{}ll@{}}
\toprule
Componente & Versión / detalle \\
\midrule
Sistema operativo & Windows 11, consola \texttt{cmd} \\
JDK & Eclipse Temurin 26.0.2 \\
Maven & Wrapper 3.9.16 (por módulo) \\
Docker Desktop & 29.9.2 \\
Python & 3.12 \\
PostgreSQL & 16, contenedor \texttt{sos-postgres}, puerto 5433 \\
MySQL & servicio local, puerto 3306 (\texttt{urgencias\_hospital}) \\
Eureka & 8761 \\
Gateway & 8080 \\
Seguridad / Pacientes / Triage & 8081 / 8082 / 8083 \\
Analítica / Agente IA & 8084 / 8088 \\
Frontend & 5500 (servidor Python del repositorio) \\
\bottomrule
\end{tabular}
\caption{Entorno usado en la fase funcional.}
\end{table}

Los servicios se levantaron con \pantalla{04. levantar-sistema.bat}. Eureka
registró los seis servicios y el gateway enroutó correctamente.

\section{Datos de prueba}

Se usó el seed canónico sin regenerarlo:

\begin{itemize}
  \item 8 perfiles, 8 actividades y 22 asignaciones.
  \item 5 categorías de triage y 3 estados.
  \item 9 cuentas aprobadas con 2FA y turno de 24 horas.
  \item 9 turnos, 24 pacientes y 30 episodios.
  \item 0 códigos 2FA, 0 tokens y 0 registros de auditoría al inicio.
\end{itemize}

Las suites crean sus propios pacientes, episodios y cuentas temporales, y los
eliminan al finalizar. Las credenciales se leen de
\pantalla{infra/seed/credenciales.local.env} y nunca se imprimen.

\section{Diseño de las suites}

\begin{table}[h]
\centering
\small
\begin{tabular}{@{}p{3.2cm}p{10.5cm}@{}}
\toprule
Suite & Qué valida \\
\midrule
\pantalla{verify_seed.py} & Login completo con 2FA de las 9 cuentas del seed, perfil en la sesión y un endpoint representativo por rol. \\
\pantalla{e2e_claves.py} & Cambio de clave, bloqueo por cambio pendiente (\texttt{ccp}), challenge 2FA de un solo uso, recuperación/activación, cambio de correo, aprobación y rechazo de cuentas, clave temporal y auditoría. \\
\pantalla{e2e_funcional.py} & Frontend, pacientes, triage, signos vitales, alta orquestada con correos y PDF, reintento, analítica, agente IA, auditoría y turnos. \\
Maven & Pruebas unitarias y de contexto de los 7 módulos (documento aparte). \\
\bottomrule
\end{tabular}
\caption{Suites ejecutadas en la fase funcional.}
\end{table}

\section{Procedimiento}

\begin{enumerate}
  \item \textbf{Preparación.} Se levantó Docker, PostgreSQL y el stack; se
  verificó el seed con \pantalla{python infra/seed/verify_seed.py}. La base
  tenía dos residuos de pruebas previas (un código 2FA y una auditoría) que se
  eliminaron para dejar el conjunto canónico.
  \item \textbf{Completar 2FA en pruebas.} Para no depender de la bandeja de
  correo, cada suite obtiene el \texttt{challenge} del login y reemplaza el
  código pendiente en la base por el hash BCrypt de \texttt{123456} usando
  \texttt{crypt()} de PostgreSQL. Después llama a
  \texttt{/auth/verificar-2fa}. Esto no modifica el código de la aplicación.
  \item \textbf{Claves y cuentas.} Se ejecutó \pantalla{e2e_claves.py}: crea
  cuentas \texttt{e2e.*} con hashes de fixture, prueba los flujos y las borra.
  \item \textbf{Flujo clínico.} Se ejecutó \pantalla{e2e_funcional.py}: crea un
  paciente y dos episodios temporales, registra dos tomas de signos, da de alta
  un episodio con la llamada única de \texttt{finalizar}, descarga el PDF,
  reenvía correos y comprueba la idempotencia del alta.
  \item \textbf{Frontend.} Se solicitan las 14 páginas y \texttt{assets/api.js}
  y se comprueba que respondan 200 con su contenido esperado.
  \item \textbf{Limpieza.} Al final se eliminan pacientes, episodios, tomas,
  auditorías y cuentas temporales. La suite admite \texttt{--no-cleanup} para
  inspección manual.
\end{enumerate}

\subsection{Consideraciones operativas}

\begin{itemize}
  \item El limitador de reenvíos 2FA permite 3 códigos por usuario cada 10
  minutos. Repetir la suite antes de esa ventana devuelve HTTP 429; se debe
  esperar o reiniciar \texttt{ms-seguridad} en local (el límite es en memoria).
  \item El alta envía correos reales a la bandeja de pruebas
  \texttt{pruebasosclinic@gmail.com}. No se usan datos clínicos reales.
  \item El agente IA consume la API de Groq; se hizo una consulta mínima.
  \item La analítica usa la base sintética \texttt{urgencias\_hospital} con
  50.000 llegadas.
\end{itemize}

\section{Criterios de aceptación}

Cada caso se registra como \ok{}, \fallo{}, \aviso{} o \omitido{}.
Un \fallo{} en una sección crítica detiene solo esa sección, no toda la suite.
Se consideran críticos la creación del paciente, la creación del episodio,
el alta persistida y la descarga del PDF. Los fallos de entrega SMTP se
registran como \aviso{} porque son una dependencia externa, siempre que el alta
se haya persistido y el reintento esté disponible.

\section{Limitaciones}

\begin{itemize}
  \item La cobertura es por API; la interfaz se valida con smoke test y una
  prueba de XSS, no con Playwright.
  \item No se midió cobertura de código (JaCoCo no está configurado).
  \item La entrega de correo y la respuesta del agente dependen de servicios
  externos (Gmail y Groq).
  \item El seed es compartido: las suites limpian sus datos, pero una
  interrupción abrupta puede dejar residuos.
\end{itemize}

\section{Reproducción}

\begin{verbatim}
04. levantar-sistema.bat
python infra/seed/verify_seed.py
python seguridad/pruebas/e2e_claves.py
python seguridad/pruebas/e2e_funcional.py
powershell -NoProfile -ExecutionPolicy Bypass `
  -File seguridad/pruebas/ejecutar-tests-maven.ps1
\end{verbatim}

\end{document}
